% section: 解法の決定

    \section{解法の決定}
      以下にここまでで扱った解法を体系的にまとめあげ,実際に問題と相対した時にどこに注目すればよいか,なにを考えればよいかをまとめた.
      チャートは3つあり,1から順に進めていく良い.
      \noindent \textbf{【チャート1】}
      \begin{enumerate}
        \item 初項の値\\
              $\rightarrow$ 不自然なら特殊関数型
        \item 何項間の漸化式か\\
              $\rightarrow$ $n$項間$(n\ge3)$なら不動点を求めて$n-1$項間に
        \item 総和を含むか\\
              $\rightarrow$ $S_{n+1}$と$S_n$の差を取って$a_{n+1}$を生み出す.
        \item 連立漸化式か\\
              $\rightarrow$ 連立漸化式なら連立漸化式型として解く
        \item 上記以外の形か\\
              $\rightarrow$ チャート2へ
      \end{enumerate}
      なお初項の値については補足を読むと良い.

      \noindent \textbf{【チャート2】}
      \begin{enumerate}
          \item 分数型か（分子分母両方に$a_n$がいるか）\\
                $\rightarrow$ $a_n$の次数が1なら分数型として解く.\,2次以上なら対数型か特殊関数型.
          \item $a_n$の係数\\
                $\rightarrow$ 定数なら特性方程式型か階差等比複合型.\,$n$に関する式の場合は階比型または対数型として解く.
          \item $a_n$の次数\\
                $\rightarrow$ 2次以上なら対数型か特殊関数型.
          \item $f(n)$が含まれている時\\
                $\rightarrow$ チャート3へ
      \end{enumerate}

      \noindent \textbf{【チャート3】}
      \begin{enumerate}
        \item $a_n$に掛けられている時\\
              $\rightarrow$ 対数型か階比型.
        \item $a_n$に足されている時\\
              $\rightarrow$ 階差型
        \item 式全体が積や商のみで構成されている時\\
              $\rightarrow$ 対数型
        \item $a_n$の項が複数存在して掛けられている時\\
              因数分解して対数型を検討.それができない場合は基本的に特殊関数型.
      \end{enumerate}
      \subsection{補足：初項の値}
        例えば初項が$\displaystyle \frac{1}{\sqrt{6}-\sqrt{2}}$だとしよう.
        わざわざこのような数字が設定されたということはきっと裏がある.
        実際これは$\cos 15^\circ$または$\sin 75^\circ$の値である.
        そこから考えて,この漸化式は三角関数型のものではないか,と予想するのである.
      
      これらに従っても解けない時は,次節を参照すること.
   
